% Compile with: uplatex file.tex && dvipdfmx file.dvi
%           or: latexmk -uplatex file.tex
\documentclass[dvipdfmx]{article}
\usepackage{amsmath, amssymb}
\usepackage{tikz}

\title{scenario\_battle\_of\_the\_sexes}
\date{\today}

\begin{document}
\maketitle

The classic two-player game Battle of the Sexes. Register player 1's (row) payoff matrix $A$ and player 2's (column) payoff matrix $B$. Row 1 / column 1 is opera, row 2 / column 2 is ballet.

$A = [2, 0; 0, 1]$

\begin{equation}
  A = \begin{pmatrix} 2 & 0 \\ 0 & 1 \end{pmatrix}
  \label{eq:E1}
\end{equation}

$B = [1, 0; 0, 2]$

\begin{equation}
  B = \begin{pmatrix} 1 & 0 \\ 0 & 2 \end{pmatrix}
  \label{eq:E2}
\end{equation}

Find the pure-strategy Nash equilibria. This scans the bimatrix for cells that are both a row and a column maximum — not expressible as a composition of general operations, so it is a genuine primitive and stays.

find the pure Nash equilibria of E1 and E2

\begin{equation}
  \left\{ \left( 1 , 1 \right) , \left( 2 , 2 \right) \right\}
  \label{eq:E3}
\end{equation}

The mixed equilibrium comes from the indifference conditions. With the opponent playing column 1 (opera) with probability $q$, player 1's rows 1 and 2 have equal expected payoff — the researcher writes that equation. Assembling it from $A$'s entries is game theory, not the system's job.

$2 q + 0 (1-q) = 0 q + 1 (1-q)$

\begin{equation}
  2 q + 0 (1-q) = 0 q + 1 (1-q)
  \label{eq:E4}
\end{equation}

Solve for $q$.

solve E4 for q

\begin{equation}
  q = \frac{1}{3}
  \label{eq:E5}
\end{equation}

Likewise, write player 2's indifference between columns from $B$'s entries.

$1 p + 0 (1-p) = 0 p + 2 (1-p)$

\begin{equation}
  1 p + 0 (1-p) = 0 p + 2 (1-p)
  \label{eq:E6}
\end{equation}

Solve for $p$. The mixed equilibrium is $(p,q)=(2/3, 1/3)$.

solve E6 for p

\begin{equation}
  p = \frac{2}{3}
  \label{eq:E7}
\end{equation}

\end{document}